\section{Current v4 Interface and Open Burdens}
\label{v480:current-v4}

This synthesis prevents the retained early material from becoming the
default status of the later v4 line. It is a reading interface, not
a replacement for the original hypotheses or finite computations.

\subsection{Shifted Normal Form and Backtranslation}

The later governing object is the shifted normal-form record of
Definition~\ref{def:v400-r05-governing-shifted-normal-form-object}.
New objects must retain occurrence class, payload/readout target,
backtranslation, guard and status. In particular packed masks must
be reconstructed from shifted supports rather than shifted as whole
integers; physical payloads must preserve their residual map and scope.
The full migration chain remains in
Appendix~\ref{v480:support-shift}. Older pre-shift wording in
Appendix~\ref{v480:support-qam} is historical input to this chain,
not authorization to undo it.

\subsection{Terminal Data and Exact Zero}

The guarded solve-chain record in
Appendix~\ref{v480:support-solve} distinguishes local C4 seam control,
finite C5 detector terminality and exact all-readout closure. The
inherited finite ledger records detector threshold $\theta=1/256$
after twenty steps; it does not state that the native residual is
exactly zero at a finite step. These recorded numbers are not newly
recovered native measurements in the present editorial pass.

The frame-typed endpoint in
Appendix~\ref{v480:support-frames} makes the distinction explicit:
matter-frame effective terminal data, a candidate photon-clock null
object and tachyon/information-facing signed interval data are
different types. None provides an automatic implication to exact
all-readout zero. A photon/clock sector is not thereby a physical
rest frame of a photon, and the terminology alone supplies no
inter-sector signal law.

\subsection{Interface Certificates and Dynamics}

The later MM/HC development passes through typed coefficients,
seed residuals, operator-domain records, local monitor certificates
and no-action/no-promotion checks. Their full dependency order is
retained in Appendix~\ref{v480:support-mm-c5}. An admissibility tag
does not define the action of a differential operator. A local
interface certificate is not a variational action or a derived
field equation.

For the proposed v4.8 metric program this leaves four concrete burdens:
declare scale units and comparison maps; realize the MM/HC operators
with their spectral/domain premises; evaluate the native residual
dictionary with fixed norms; and supply a dynamical realization with
its approximation and stability bounds. The earlier QAM/router
construction can record these objects but cannot replace them.

Section~\ref{v480:hc-dynamics-geometry} now supplies explicit conditional
calculations rather than a dynamical identification: projected hidden
states and memory, conservative collective inertia, the null HC kinetic
tensor, real-action/current alternatives and an off-shell geometry/stress
match. Proposition~\ref{v480:hc-same-frame-obstruction} excludes the direct
same-induced-frame real-bilinear master/coframe identification under its
stated assumptions. A separate microscopic/effective geometry remains
possible but needs its own common action and current balance.

Section~\ref{v490:cartan-source} now supplies a declared neutral Cartan
comparison with a finite normal inverse and a conditional source law.
Preserving compatible tangential jets determines the six source values
at the studied derivative order; unrestricted higher jets can instead
absorb them. Neither fact establishes arbitrary surface admissibility,
the full regional equations or the native MM/HC identification.
